% \iffalse meta-comment
% LTeX: language=de-DE
%
%  TUDCD-Script -- Corporate Design of Technische Universität Dresden
% ----------------------------------------------------------------------------
%
%  Copyright (C) TU Dresden <tudlatex@tu-dresden.de>, 2025
%
% ----------------------------------------------------------------------------
%
%  This work may be distributed and/or modified under the conditions of the
%  LaTeX Project Public License, either version 1.3c of this license or
%  any later version. The latest version of this license is in
%    http://www.latex-project.org/lppl.txt
%  and version 1.3c or later is part of all distributions of
%  LaTeX version 2008-05-04 or later.
%
%  This work has the LPPL maintenance status "maintained".
%
%  The current maintainer and author of this work is Yoshi Diepelt.
%
% ----------------------------------------------------------------------------
% \fi
% \selectlanguage{ngerman}
%
% \subsection{\texttt{test-auto-options}: Tests für die Optionenauswahl}
%
% Um bei Regressionstests neue Optionen sicher einführen zu können,
% werden die einzelnen Optionen der entsprechenden Klassen auf korrekte Funktionsweise überprüft.
%
% Dabei werden für die einzelnen Optionstypen \emph{global},\emph{spät} und \emph{einfach}
% folgende Tests durchgeführt: \begin{description}
%   \item[Globale Optionen]
%   Für globale Optionen ist es notwendig, dass diese nach der möglichen Angabe in \dmacro{\documentclass}
%   ungültig sind, und den Übersetzungsprozess abbrechen, sollte detektiert werden,
%   dass sie dennoch geladen werden.
%   \item[Späte Optionen]
%   Für späte Optionen ist es notwendig, dass diese nach der Dokumentenpräambel im Dokument ungültig sind.
%   \item[Einfache Optionen]
%   Für einfache Optionen ist es notwendig, dass diese jederzeit geändert werden können, und die Änderung
%   nicht zu Übersetzungsabrüchen führt.
% \end{description}
%
% Beim Entpacken werden die Testdokumente innerhalb einer \dcode{\textbackslash ifx}-Abfrage übersprungen.
% Damit die darin genutzten Schalter von \LaTeX{} korrekt verschachtelt werden, müssen sie beim Entpacken
% als Bedingungen bekannt sein.
%    \begin{macrocode}
%<*ins>
\input docstrip.tex
\def\tudcdbatchfilemode{\relax}
\let\if@altsecnumformat\iffalse
\let\if@openright\iffalse
%</ins>
%<*ins>
\ifx\tudcdbatchfilemode\undefined%
%</ins>
%<*preamble>
\input{regression-test.tex}
%</preamble>
%<*ins>
\fi
%</ins>
%    \end{macrocode}
%
% \subsubsection{Testen der Optioneninferenz}
%
% \begin{unittest}{dev-version}{Automatische Inferenz der Logosprache: \texttt{logo/language}}
%
% In dem ersten Test wird eine erfolgreiche Kompilation der Artikelklasse
% getestet.
%
%    \begin{macrocode}
%<*ins>
\begingroup
\def\testsuiteglobalforclass#1{
  \file{\jobname-0001-logolang-#1-global-babel-en.lvt}{\from{\jobname.dtx}{0001,preamble,#1,babel,global,english,test}}
  \file{\jobname-0001-logolang-#1-global-babel-en.lve}{\from{\jobname.dtx}{0001,preamble,#1,babel,global,english}}
  %%
  \file{\jobname-0001-logolang-#1-global-babel-de.lvt}{\from{\jobname.dtx}{0001,preamble,#1,babel,global,german,test}}
  \file{\jobname-0001-logolang-#1-global-babel-de.lve}{\from{\jobname.dtx}{0001,preamble,#1,babel,global,german}}
}

\def\testsuitelateforclass#1{
  \file{\jobname-0001-logolang-#1-late-babel-en.lvt}{\from{\jobname.dtx}{0001,preamble,#1,babel,late,english,test}}
  \file{\jobname-0001-logolang-#1-late-babel-en.lve}{\from{\jobname.dtx}{0001,preamble,#1,babel,late,english}}
  %%
  \file{\jobname-0001-logolang-#1-late-babel-de.lvt}{\from{\jobname.dtx}{0001,preamble,#1,babel,late,german,test}}
  \file{\jobname-0001-logolang-#1-late-babel-de.lve}{\from{\jobname.dtx}{0001,preamble,#1,babel,late,german}}
}
\generate{
  \testsuiteglobalforclass{article}
  \testsuitelateforclass{article}
  \testsuiteglobalforclass{report}
  \testsuitelateforclass{report}
  \testsuiteglobalforclass{book}
  \testsuitelateforclass{book}
  \testsuiteglobalforclass{thesis}
  \testsuitelateforclass{thesis}
  \testsuiteglobalforclass{letter}
  \testsuitelateforclass{letter}
  \testsuitelateforclass{beamer}
}
\endgroup
%</ins>
%    \end{macrocode}
% Und dies ohne jegliche zusätzliche Pakete.
%    \begin{macrocode}
%<*ins>
\ifx\tudcdbatchfilemode\undefined%
%</ins>
%<*0001>
%<@@=tudcd>
\documentclass
%<global>[logo/language=auto]
{
%<article>tudcdartcl%%
%<report>tudcdreprt%%
%<book>tudcdbook%%
%<thesis>tudcdthesis%%
%<beamer>beamer%%
%<letter>tudcdlttr%%
}
%<beamer>\usetheme{tudcd}
%<babel&german>\usepackage[ngerman]{babel}
%<babel&english>\usepackage[english]{babel}
%<*late>
\TUDCDsetup{
  logo/language=auto
}
%</late>

\begin{document}
\showoutput%
\vfil\break%
\START
%<*test>
\ExplSyntaxOn%
\typeout{\str_use:N\l_@@_logolang_str}
\ExplSyntaxOff%
%</test>
%<!test&german>\typeout{de}
%<!test&english>\typeout{en}
\END
\end{document}
%</0001>
%<*ins>
\fi
%</ins>
%    \end{macrocode}
%
% \end{unittest}
%
% \subsubsection{Testen der Weitergabe an \KOMAScript}
%
% \begin{unittest}{dev-version}{Weiterreichen unbekannter Optionen an die \KOMAScript-Klasse}
%
% Optionen, welche die Klassen nicht selbst verarbeiten, werden an die zugrundeliegende
% \KOMAScript-Klasse weitergereicht.
% Dabei müssen sowohl Optionen der Form \dcode{<Schlüssel>=<Wert>} als auch Optionen
% ohne Wert vollständig ankommen.
% Als erwartetes Ergebnis dient die jeweilige \KOMAScript-Klasse mit denselben Optionen.
%
%    \begin{macrocode}
%<*ins>
\begingroup
\def\testsuitekomaforclass#1{
  \file{\jobname-0002-komaoptions-#1.lvt}{\from{\jobname.dtx}{0002,preamble,#1,test}}
  \file{\jobname-0002-komaoptions-#1.lve}{\from{\jobname.dtx}{0002,preamble,#1}}
}
\generate{
  \testsuitekomaforclass{article}
  \testsuitekomaforclass{report}
  \testsuitekomaforclass{book}
  \testsuitekomaforclass{thesis}
}
\endgroup
%</ins>
%    \end{macrocode}
% Die Optionen \doption{numbers}, \doption{parskip} und \doption{listof} werden mit Wert,
% \doption{openright} ohne Wert angegeben.
% Ausgegeben wird deren Auswirkung auf die interne Konfiguration von \KOMAScript.
%    \begin{macrocode}
%<*ins>
\ifx\tudcdbatchfilemode\undefined%
%</ins>
%<*0002>
\documentclass[numbers=enddot,parskip=half,listof=totoc,openright]
{
%<article&test>tudcdartcl%%
%<report&test>tudcdreprt%%
%<book&test>tudcdbook%%
%<thesis&test>tudcdthesis%%
%<article&!test>scrartcl%%
%<report&!test>scrreprt%%
%<book&!test>scrbook%%
%<thesis&!test>scrbook%%
}

\begin{document}
\START
\makeatletter
\typeout{numbers=\if@altsecnumformat enddot\else noenddot\fi}
\typeout{parskip=\ifdim\parskip>\z@ half\else false\fi}
\typeout{listof=\Iftocfeature{lof}{totoc}{totoc}{nototoc}}
\typeout{openright=\@ifundefined{if@openright}{-}{\if@openright true\else false\fi}}
\makeatother
\END
\end{document}
%</0002>
%<*ins>
\fi
%</ins>
%    \end{macrocode}
%
% \end{unittest}
% \subsubsection{Testen des Ausgabemediums}
%
% \begin{unittest}{dev-version}{Voreinstellungen der Option \texttt{output} der Klasse für Abschlussarbeiten}
%
% Die Option \doption{output} setzt Voreinstellungen für \KOMAScript{} sowie die Farben.
% Explizit gesetzte Optionen müssen dabei Vorrang haben.
% Ausgegeben werden die doppelseitige Einstellung, \doption{open}, die Bindekorrektur in Millimetern,
% die Hervorhebungsfarbe und die Farbe des Logos.
%    \begin{macrocode}
%<*ins>
\begingroup
\def\testsuiteoutput#1{
  \file{\jobname-0003-output-#1.lvt}{\from{\jobname.dtx}{0003,preamble,#1,test}}
  \file{\jobname-0003-output-#1.lve}{\from{\jobname.dtx}{0003,preamble,#1}}
}
\generate{
  \testsuiteoutput{digital}
  \testsuiteoutput{print}
  \testsuiteoutput{printoneside}
  \testsuiteoutput{printbcor}
  \testsuiteoutput{digitaltwoside}
  \testsuiteoutput{printcolor}
}
\endgroup
%</ins>
%    \end{macrocode}
%
%    \begin{macrocode}
%<*ins>
\ifx\tudcdbatchfilemode\undefined%
%</ins>
%<*0003>
%<@@=tudcd>
\documentclass
%<digital>[output=digital]
%<print>[output=print]
%<printoneside>[output=print,oneside]
%<printbcor>[output=print,BCOR=5mm]
%<digitaltwoside>[twoside]
%<printcolor>[output=print,color/highlight-color=Rot1,logo/color=blue]
%<test>{tudcdthesis}
%<!test>{article}

\begin{document}
\START
%<*test>
\makeatletter
\ExplSyntaxOn
\typeout{twoside=\legacy_if:nTF{@twoside}{true}{false}}
\typeout{openright=\legacy_if:nTF{@openright}{true}{false}}
\typeout{BCOR=\fp_eval:n{round(\dim_to_decimal_in_unit:nn{\ta@bcor}{1mm},1)}}
\typeout{color=\tl_use:N\l_@@_highlight_color_tl}
\typeout{logo=\str_use:N\l_@@_logocolor_str}
\ExplSyntaxOff
\makeatother
%</test>
%<*!test>
%<digital|digitaltwoside>\typeout{twoside=%
%<digital>false}
%<digitaltwoside>true}
%<print|printbcor|printcolor>\typeout{twoside=true}
%<printoneside>\typeout{twoside=false}
\typeout{openright=true}
%<digital|digitaltwoside|printoneside|print|printcolor>\typeout{BCOR=%
%<digital|digitaltwoside>0}
%<print|printoneside|printcolor>10}
%<printbcor>\typeout{BCOR=5}
%<digital|digitaltwoside>\typeout{color=Brilliantblau}
%<print|printoneside|printbcor>\typeout{color=Schwarz}
%<printcolor>\typeout{color=Rot1}
%<digital|digitaltwoside|printcolor>\typeout{logo=blue}
%<print|printoneside|printbcor>\typeout{logo=black}
%</!test>
\END
\end{document}
%</0003>
%<*ins>
\fi
%</ins>
%    \end{macrocode}
%
% \end{unittest}
%    \begin{macrocode}
%<*ins>
\endbatchfile
%</ins>
%    \end{macrocode}
